% ------------------------------------------------------------------
% TWIN ARITHMETIC & MONOGENIC GROUPS — CLEAN MASTER TEMPLATE
% ------------------------------------------------------------------

\documentclass[11pt]{article}

% ----------------------------
% PACKAGES
% ----------------------------
\usepackage{amsmath,amssymb,amsfonts,amsthm}
\usepackage{mathrsfs}
\usepackage{enumitem}
\usepackage{hyperref}
\usepackage{geometry}
\usepackage{mathtools}
\usepackage{stmaryrd}
\usepackage{microtype}

\geometry{margin=1in}

% ----------------------------
% THEOREM ENVIRONMENTS
% ----------------------------
\newtheorem{theorem}{Theorem}[section]
\newtheorem{lemma}[theorem]{Lemma}
\newtheorem{proposition}[theorem]{Proposition}
\newtheorem{corollary}[theorem]{Corollary}

\theoremstyle{definition}
\newtheorem{definition}[theorem]{Definition}

\theoremstyle{remark}
\newtheorem{remark}[theorem]{Remark}

% ----------------------------
% MACROS
% ----------------------------
\newcommand{\Z}{\mathbb{Z}}
\newcommand{\R}{\mathbb{R}}
\newcommand{\N}{\mathbb{N}}
\newcommand{\Q}{\mathbb{Q}}
\newcommand{\C}{\mathbb{C}}

\newcommand{\opn}{\odot_n}
\newcommand{\opzero}{\odot_0}
\newcommand{\phin}{\phi^n}
\newcommand{\phim}{\phi^{-n}}

% ------------------------------------------------------------------
% BEGIN DOCUMENT
% ------------------------------------------------------------------

\begin{document}
	
	\begin{center}
		{\LARGE \textbf{Twin Arithmetic and Monogenic Transport of Ring Structures}}\\[1em]
		{\large Jocelyn Nembe}
	\end{center}
	
	\vspace{1.5em}
	
	\begin{abstract}
		We develop a general theory of arithmetic hierarchies obtained by transporting
		the algebraic structure of a ring through the action of a monogenic group of
		applications.  
		Given a bijection $\phi:A\to A$ on a ring $(A,+,\cdot)$,
		the iterates $\phi^n$ induce a family of conjugate operations 
		$(\oplus_n,\otimes_n)$ forming a hierarchy of ``twin rings'',
		all isomorphic to the base ring.  
		This indexing by the additive group $(\Z,+)$ leads to a canonical theory of
		prime operations, factorization of operations, base-change identities,
		and arithmetical decompositions mirroring the classical decomposition of
		integers into prime factors.  
		The exponential case $\phi=\exp$ recovers ordinary addition and multiplication
		as two consecutive levels of the same structure, while higher levels
		produce non-linear arithmetic systems.  
		This article establishes the algebraic foundations of these twin hierarchies
		and prepares the study of growth laws, combinatorics, and higher-order
		arithmetics.
	\end{abstract}
	
	% ================================================================
	% SECTION 1 — INTRODUCTION
	% ================================================================
	\section{Introduction}
	
	In the classical approach, a binary operation 
	\[
	\odot : E\times E \to E
	\]
	is a fixed algebraic object defined in a single coordinate system.
	In the framework of \emph{Twin Spaces}, mathematical structures are no longer
	studied in isolation but as families indexed by a symmetry group acting on
	the underlying set.  
	Given a group $G$ acting on $E$, the fundamental compatibility rule for an
	equivariant operation reads
	\[
	(x \odot_{e,g} y)
	= \overline{g}\,\bigl(x\odot_e y\bigr),
	\]
	where $e$ is the identity of $G$, and $\overline{g}=g^{-1}$.
	
	In this article we focus on the case where the symmetry group is 
	\emph{monogenic}, generated by a bijection $\phi:E\to E$.
	Thus,
	\[
	G = \langle \phi\rangle = \{\phi^n : n\in\Z\}
	\]
	with composition as group law.  
	This group is always isomorphic to $(\Z,+)$, and its index $n$ acts as an
	``arithmetical coordinate'' used to generate a hierarchy of operations
	and, more generally, entire algebraic structures.
	
	\paragraph{Transported arithmetic.}
	Given a reference operation $\odot_0$ on $E$, the twin transport rule produces,
	for each integer $n$,
	\[
	x \odot_n y
	=
	\phi^{n}\bigl(\phi^{-n}(x)\odot_0\phi^{-n}(y)\bigr).
	\]
	The family $\{\odot_n\}_{n\in\Z}$ forms a hierarchy of conjugate operations.
	Every level is algebraically equivalent to the reference level via
	$\phi^{-n}$, and the hierarchy inherits the additive structure of $\Z$.
	
	\paragraph{Extension to rings.}
	The construction extends verbatim to any ring $(A,+,\cdot)$:
	\[
	x\oplus_n y = \phi^{n}\bigl(\phi^{-n}(x)+\phi^{-n}(y)\bigr),\qquad
	x\otimes_n y = \phi^{n}\bigl(\phi^{-n}(x)\cdot\phi^{-n}(y)\bigr),
	\]
	producing twin rings $(A,\oplus_n,\otimes_n)$ all isomorphic to $(A,+,\cdot)$.
	
	\paragraph{Base-change and primality.}
	The group isomorphism $G\simeq\Z$ allows transportation of classical arithmetic
	to the level of operations:
	\[
	x_1 \odot_{n+k} x_2
	=
	\phi^{k}\bigl(\phi^{-k}(x_1)\odot_n \phi^{-k}(x_2)\bigr).
	\]
The integer index also carries the \emph{multiplicative} structure of $\N^*$: by the
iteration identity $\phi^{ab}=(\phi^a)^{\circ b}$, a factorization $q=p_1^{\alpha_1}\cdots p_r^{\alpha_r}$
realizes $\odot_q$ as a nesting of the prime-level constructions $\odot_{p_i}$. This is unique
factorization in $(\N^*,\times)$ --- a structure on the index distinct from the additive group
law $\phi^n\phi^m=\phi^{n+m}$ of $G\simeq(\Z,+)$.
	
	\paragraph{The exponential case.}
	When $\phi=\exp$ (on any invertible domain), the twin hierarchy recovers:
	\[
	x\odot_0 y=x+y,\qquad x\odot_1 y = xy,
	\]
	and higher levels yield non-linear ``exponential'' arithmetics.
	Addition and multiplication appear as two consecutive states in the same 
	transported structure.
	
	\subsection*{Historical and conceptual motivation}
	
	The idea of iterating operations is as old as classical hyper-operations
	and the Ackermann hierarchy.  
	However, these constructions are usually defined syntactically rather
	than algebraically and depend on specific choices of analytic functions.
	
	The twin approach is different in nature:
	\begin{itemize}[leftmargin=2em]
		\item It applies to \emph{any} ring, not only to $\R$ or to analytic settings.
		\item It is driven by \emph{symmetry} rather than by combinatorial iteration.
		\item It produces a \emph{family of isomorphic algebraic structures}, not a
		single growing function.
		\item It yields a true \emph{arithmetic of operations}, including primality,
		factorization, and equivalence classes.
	\end{itemize}
	
	This makes twin arithmetic both a generalization of classical operations 
	and a new framework for understanding algebraic transformations through 
	group actions.
	
	\vspace{1em}
	
	The goal of this article is to lay the algebraic foundations of
	twin arithmetic and prepare for the study of higher-level structures,
	including growth laws, combinatorial identities, and applications.
	
	% ================================================================
	% SECTION 2 — TWIN STRUCTURES ON RINGS
	% ================================================================
	\section{Twin Structures on Rings}
	\label{sec:2}
	
	\subsection{Monogenic actions and transported ring operations}
	
	Let $(A,+,\cdot)$ be any ring (commutative or not).
	Let $\phi:A\to A$ be a bijection.
	We define the monogenic group
	\[
	G=\langle \phi\rangle=\{\phi^n : n\in\Z\}
	\]
	acting on $A$ by:
	\[
	\phi^n\cdot x = \phi^n(x).
	\]
	
	\begin{definition}[Twin ring operations]
		For each $n\in\Z$ define
		\[
		x\oplus_n y := \phi^n\bigl(\phi^{-n}(x)+\phi^{-n}(y)\bigr),
		\qquad
		x\otimes_n y := \phi^n\bigl(\phi^{-n}(x)\cdot\phi^{-n}(y)\bigr).
		\]
	\end{definition}
	
	\begin{proposition}[Twin rings]
		Each $(A,\oplus_n,\otimes_n)$ is a ring.
		Moreover,
		\[
		\Phi_n:A\to A,\quad \Phi_n(x)=\phi^{-n}(x)
		\]
		yields an isomorphism
		\[
		\Phi_n: (A,\oplus_n,\otimes_n)\xrightarrow{\ \sim\ }(A,+,\cdot).
		\]
	\end{proposition}
	\begin{proof}
		The map $\Phi_n=\phi^{-n}$ is a bijection of $A$, and by definition
		$\Phi_n(x\oplus_n y)=\phi^{-n}\phi^{n}\!\bigl(\phi^{-n}(x)+\phi^{-n}(y)\bigr)=\Phi_n(x)+\Phi_n(y)$,
		and likewise $\Phi_n(x\otimes_n y)=\Phi_n(x)\cdot\Phi_n(y)$. Thus $\Phi_n$ carries
		$(A,\oplus_n,\otimes_n)$ bijectively onto $(A,+,\cdot)$ while preserving both operations;
		transporting the ring axioms back along $\Phi_n^{-1}=\phi^{n}$ shows $(A,\oplus_n,\otimes_n)$
		is a ring, with zero $\phi^{n}(0)$, unit $\phi^{n}(1)$ and additive inverse
		$\ominus_n x=\phi^{n}(-\phi^{-n}x)$. Hence every twin level is a ring, isomorphic to the base
		ring via $\Phi_n$.
	\end{proof}
	
	\subsection{Base-change identity and factorization}
	
	\begin{proposition}[Base-change identity]
		For all $n,k\in\Z$,
		\[
		x_1\odot_{n+k}x_2
		=
		\phi^{k}\bigl(\phi^{-k}(x_1)\odot_n \phi^{-k}(x_2)\bigr).
		\]
	\end{proposition}
	\begin{proof}
		Unwinding the right-hand side through the defining formula,
		\[
			\phi^{k}\!\bigl(\phi^{-k}(x_1)\odot_n\phi^{-k}(x_2)\bigr)
			=\phi^{k}\phi^{n}\!\bigl(\phi^{-n}\phi^{-k}(x_1)\odot_0\phi^{-n}\phi^{-k}(x_2)\bigr)
			=\phi^{n+k}\!\bigl(\phi^{-(n+k)}(x_1)\odot_0\phi^{-(n+k)}(x_2)\bigr),
		\]
		which is $x_1\odot_{n+k}x_2$, using $\phi^{k}\phi^{n}=\phi^{n+k}$ and
		$\phi^{-n}\phi^{-k}=\phi^{-(n+k)}$.
	\end{proof}
	
The group law of $G\simeq(\Z,+)$ is additive: composition shifts the level,
$\phi^n\phi^m=\phi^{n+m}$, and the base-change identity expresses this. A separate,
\emph{multiplicative} structure on the index arises from iteration, $\phi^{ab}=(\phi^a)^{\circ b}$;
with respect to it the levels inherit the prime/composite behavior of $(\N^*,\times)$.
	
	\begin{definition}[Prime and composite operations]
		For $n\ge 2$ we call $\odot_n$ \emph{prime} if the integer $n$ is prime --- equivalently, if
		$\phi^n$ is not a non-trivial iterate $(\phi^d)^{\circ(n/d)}$ for any $1<d<n$ --- and
		\emph{composite} otherwise. This refers to the multiplicative structure of the index
		$n\in\N^*$, not to the additive group law of $G$.
	\end{definition}
	
	\begin{proposition}[Iterated structure of composite levels]
		Let $q=p_1^{\alpha_1}\cdots p_r^{\alpha_r}$ be the prime factorization of $q$ in
		$(\N^*,\times)$. By the iteration identity $\phi^{ab}=(\phi^a)^{\circ b}$, the operation
		$\odot_q$ is the corresponding nesting of the prime-level twin constructions: writing
		$\psi=\phi^{a}$, the level-$ab$ operation $\odot_{ab}$ equals the level-$b$ twin built from
		$\psi$. This nesting is unique up to the order of the (commuting) factors, the uniqueness
		being that of the integer factorization of $q$. It is \emph{not} a decomposition under the
		additive group law $\phi^n\phi^m=\phi^{n+m}$ of $G\simeq(\Z,+)$.
	\end{proposition}
	\begin{proof}
		Write $\psi=\phi^{a}$, so that $\psi^{\circ b}=(\phi^{a})^{\circ b}=\phi^{ab}$. The level-$b$ twin
		operation built from the bijection $\psi$ is
		$x\mapsto\psi^{b}\!\bigl(\psi^{-b}(x)\odot_0\psi^{-b}(y)\bigr)
		=\phi^{ab}\!\bigl(\phi^{-ab}(x)\odot_0\phi^{-ab}(y)\bigr)=x\odot_{ab}y$. Applying this with the
		successive prime-power factors of $q$ realizes $\odot_q$ as the iterated nesting of the
		prime-level constructions. The factors commute and the factorization of $q$ is unique
		(fundamental theorem of arithmetic), so the nesting is unique up to the order of factors.
	\end{proof}

	\begin{remark}
		The two structures on the index must not be conflated. As a \emph{group},
		$G=\langle\phi\rangle\simeq(\Z,+)$ is generated by $\phi$ and has no notion of primality;
		the base-change identity expresses the additive (torsor) structure of $\{\odot_n\}$.
		Primality and factorization live in the \emph{multiplicative} monoid $(\N^*,\times)$ and are
		realized on operations through iteration of the transport, $\phi^{ab}=(\phi^a)^{\circ b}$.
	\end{remark}
	
	% ================================================================
	% SECTION 3 — EXPONENTIAL AND \phi-HIERARCHIES
	% ================================================================
	\section{Hierarchies from Monogenic Transport}
	\label{sec:3}
	
	Let $(A,+,\cdot)$ be a ring, and $\phi:A\to A$ a bijection.
	We obtain the transported operations
	\[
	x\odot_n y = \phi^{n}\bigl(\phi^{-n}(x)\odot_0\phi^{-n}(y)\bigr).
	\]
	
	\subsection{The exponential case}
	
	When $\phi=\exp$ on a suitable domain,
	\[
	x\odot_0 y=x+y,\qquad
	x\odot_1 y=xy,\qquad
	x\odot_2 y=\exp(\log(x)\log(y)).
	\]
	Higher levels produce non-linear arithmetic systems 
	with structurally inherited properties.

	\begin{remark}[Domains in the exponential case]
		On $\R$ the map $\exp$ is a bijection only onto $\R_{>0}$, so the iterates $\phi^{\pm n}$, and
		hence $\odot_n$, are defined only on the nested domains where the relevant logarithms exist; for
		example $\odot_2(x,y)=\exp(\log x\,\log y)=x^{\log y}$ requires $x,y>0$. As with all iterated
		$\exp/\log$ towers, each higher level shrinks the natural domain. On a ring where $\phi$ is a
		genuine bijection (a formal, $p$-adic, or finite setting) no such restriction arises and the whole
		hierarchy $\{\odot_n\}_{n\in\Z}$ is globally defined.
	\end{remark}
	
	\subsection{Aggregated operations}
	
	If $\odot_0$ is associative,
	\[
	\bigodot_{i=1}^k{}^{(n)} x_i
	=
	\phi^n\left(\bigodot_{i=1}^k{}^{(0)}\phi^{-n}(x_i)\right),
	\]
	generalizing exponential aggregation, log-domain linearization, and 
	other classical transforms.
	
	% ================================================================
	% PLACEHOLDER FOR NEXT SECTIONS
	% ================================================================
	\section{Growth Laws and Higher-order Operations}

	\begin{proposition}[Order preservation]
		Suppose $A$ is totally ordered and $\phi$ is strictly increasing. Then for every $n$ the twin
		operation $\odot_n$ is strictly increasing in each argument, and
		$\phi^{-n}:(A,\odot_n)\to(A,\odot_0)$ is an order isomorphism. In particular all levels share the
		monotonicity type of the base operation.
	\end{proposition}
	\begin{proof}
		The maps $\phi^{n}$ and $\phi^{-n}$ are strictly increasing (composites of a strictly increasing
		bijection), and $x\odot_n y=\phi^{n}(\phi^{-n}x\odot_0\phi^{-n}y)$ is a composite of increasing maps
		with the increasing $\odot_0$, so it is strictly increasing in each argument. The displayed map is
		an order isomorphism by the Twin-rings proposition together with this monotonicity.
	\end{proof}

	\begin{remark}[The neutral elements move with the level]
		The level-$n$ ring has zero $0_n=\phi^{n}(0)$, unit $1_n=\phi^{n}(1)$, additive inverse
		$\ominus_n x=\phi^{n}(-\phi^{-n}x)$, and, where defined, multiplicative inverse
		$\phi^{n}\!\bigl((\phi^{-n}x)^{-1}\bigr)$. The neutral elements therefore translate with $n$: for
		$\phi=\exp$ the additive zero $0_0=0$ becomes the multiplicative unit $1_1=e^{0}=1$ at the next
		level, which is precisely why level-$0$ addition becomes level-$1$ multiplication.
	\end{remark}

	\begin{remark}[Twin hierarchy versus hyperoperations]
		For $\phi=\exp$ the sequence $\odot_0=+$, $\odot_1=\times$ continues with
		$\odot_2(x,y)=x^{\log y}$, a commutative \emph{associative} third operation conjugate to $+$. This
		is \emph{not} the Ackermann/Goodstein hyperoperation sequence, whose third term (tetration) is
		already neither associative nor commutative. Every twin level remains a ring operation isomorphic
		to the base; the two sequences agree only at levels $0$ and $1$. The ``growth'' of a level is thus
		entirely carried by the stretching of $\phi^{n}$, not by a change in algebraic type.
	\end{remark}
	
	\section{Prime Decompositions and Twin Complexity}

	\begin{definition}[Twin complexity]
		For $q\ge 1$ define the \emph{twin complexity} $\Omega(\odot_q):=\Omega(q)$, the number of prime
		factors of $q$ counted with multiplicity, with $\Omega(\odot_1)=0$.
	\end{definition}

	\begin{proposition}[Minimal nesting depth]
		The minimal number of prime-level twin constructions whose iterated nesting realizes $\odot_q$
		equals $\Omega(q)$.
	\end{proposition}
	\begin{proof}
		By the proposition on iterated structure, each prime factor $p$ of $q$ contributes exactly one
		level-$p$ nesting, so a factorization of $q$ into $k$ primes (with multiplicity) yields a nesting of
		depth $k$. The fundamental theorem of arithmetic gives $k=\Omega(q)$. For minimality, $m$ prime
		constructions multiply the index by a product of $m$ primes, hence reach an index with at most $m$
		prime factors; reaching $q$ therefore requires $m\ge\Omega(q)$.
	\end{proof}

	\begin{remark}[Additive cost versus twin complexity]
		Twin complexity measures cost in the multiplicative monoid $(\N^*,\times)$ --- the nesting depth of
		transports. The additive group $G\simeq(\Z,+)$ carries a different cost: the word length $|n|$ of
		$\phi^{n}$ in the generator $\phi$. The two are unrelated. For instance $\odot_{2^{10}}$ has additive
		cost $1024$ but twin complexity $10$, whereas $\odot_p$ for a large prime $p$ has additive cost $p$
		but twin complexity $1$.
	\end{remark}
	
	\section{Applications}

	The twin hierarchy exhibits several familiar constructions as instances of a single transport.

	\paragraph{Log-domain arithmetic and numerical stability.}
	For $\phi=\exp$, level $1$ is multiplication and level $0$ is addition, so evaluating a product in the
	$\phi^{-1}=\log$ coordinate replaces multiplication by addition --- the standard log-domain trick. The
	level-$(-1)$ operation $x\odot_{-1}y=\log(e^{x}+e^{y})$ is the \emph{LogSumExp} function, the
	numerically stable evaluation of a sum of exponentials ubiquitous in statistics and machine learning.
	Both are twin operations of the single bijection $\exp$.

	\paragraph{Relativistic addition.}
	For $\phi=\operatorname{artanh}(\,\cdot\,/c)$ on $(-c,c)$ with $\odot_0=+$, level $1$ is Einstein
	velocity addition $x\odot_1 y=(x+y)/(1+xy/c^{2})$, and the hierarchy iterates the rapidity map. The
	associativity of velocity addition is inherited, through $\phi$, from associativity of $+$.

	\paragraph{$q$-arithmetic.}
	Taking $\phi$ a $q$-exponential recovers the $q$-addition and $q$-multiplication of nonextensive
	statistics~\cite{tsallis2009} as adjacent twin levels, so that their algebraic identities are
	transported instances of the ordinary ring axioms.

	In each case the value of the twin formalism is uniformity: a single bijection $\phi$ and its integer
	iterates organize, as one hierarchy, operations that are classically studied in isolation.
	
	\begin{thebibliography}{9}
		\bibitem{nembe-foundations} J.~Nembe, \emph{Twin Spaces: Foundations}, ROOTS-INSIGHTS Vol.~1 (2026).
		\bibitem{hardy-wright} G.~H.~Hardy and E.~M.~Wright, \emph{An Introduction to the Theory of Numbers}, 6th ed., Oxford University Press, 2008.
		\bibitem{tsallis2009} C.~Tsallis, \emph{Introduction to Nonextensive Statistical Mechanics}, Springer, 2009.
	\end{thebibliography}

\end{document}
